\cotefile{1}
\scanpage{4}
Soit $C$ une $C^*$-algèbre, on veut déterminer les \deleted{\unsure{endomorphismes}} endomorphismes linéaires continus $T$ permutant aux translations gauches. Soit $T$ un \unsure{telle|tel}, son bitransposé $T''$ est un endomorphisme \unsure{fblt} continu \inserted{de $C''$} qui permute aux translations $x\unsure{\mapsto|\to}xy$ \deleted{$(y\in C)$ donc} pour $y\in C$, donc pour \deleted{tout} $y\in C''$ (\inserted{donc} $T''(xy)=T''(x)y$ \deleted{\unsure{il est} \illegible{2} \unsure{donc}} \inserted{les deux} mbs sont \unsure{fblt} continus). Donc de la forme $x\to ux$ ($u\in C''$ \inserted{bien déterminé de \unsure{norme} $\unsure{\lVert T''\rVert|\lVert T\rVert}$}). En fait, $u$ appartient à l'ens. $\Gamma_g$ des éléments $u$ de $C''$ tels que $uC\subset C$. On voit donc que l'ensemble d'endomorphismes cherché est isomorphe, avec sa structure d'algèbre normée, à $\Gamma_g$, \deleted{\unsure{l'hermitien}} $u$ hermitien signifie $x^*T(x)$ hermitien pour \unsure{tout} $x\in C$ ; $u$ positif signifie $x^*T(x)\geq 0$ pour \deleted{tout} $\unsure{u|x}\in C$.

Soit \deleted{$\unsure{\mathfrak g|g}$} $\Gamma$ l'ens des $u\in C''$ tels que $u\underline{C}\subset\underline{C}$, $\underline{C}u\subset\underline{C}$. C'est une sous $C^*$-algèbre \unsure{involutive}. \unsure{Tout} \unsure{est|elt} \unsure{alors|élément} hermitien de $\Gamma_g$ est dans $\Gamma$

Supposons maintenant $C$ \unsure{plongé|prolongé} dans la $C^*$-algèbre $R$, et soit $\Gamma_1$ l'ens des $x\in R$ tels que $xC\subset C$, $Cx\subset \unsure{R|C}$, c'est une sous $C^*$-algèbre de $R$. Pour tout $a\in\Gamma_1$, soit $\varphi(a)$ l'élément de $C''$ correspondant à l'opération $T_a(x)=ax$ dans $C$. Si $a$ hermitien, $\varphi(a)$ est hermitien, donc dans $\Gamma$, donc pour \deleted{tout} $a\in\Gamma_{\unsure{g|1}}$, $\varphi(a)\in\Gamma$. De plus, $\varphi$ est évidemment une représentation d'algèbre, compatible avec l'involution (puisque
% archive sheet number: 5

\scanpage{12}
2. \uline{Fonctions de Weyl}

% slanted
\margin{$\unsure{K^\infty|L^\infty}, L^{\unsure{\infty(K)|\infty}}_{\unsure{K|k}}$ \unsure{changer|toujours} \unsure{$L$} \unsure{en|sur} $K^\infty_{\unsure{K|k}}$} \uline{Définition 1} \deleted{Soit} \deleted{\unsure{$M$} \unsure{un mesuré} (\unsure{mesure} \unsure{$m$}).} On appelle \uline{fonction de Polya} sur \deleted{\unsure{$M$}} \inserted{$I=]0,\omega[$ $(0<\omega\leq+\infty)$} une fonction $P$ réelle définie sur l'espace $K^\infty(I)$ des fonctions mesurables à support compact, ayant les propriétés suivantes :

$P_1\unsure{)|/}$ $P$ est \uline{symétrique}, i.e. $P(f)=P(g)$ si $f$ et $g$ sont spectralement équivalentes

$P_2)$ $P$ est convexe, i.e. $P(\lambda f+(1-\lambda)g)\leq\lambda P(f)+(1-\lambda)P(g)$ si $0\leq\lambda\leq1$

$P_3)$ \unsure{Pour|Soit} \deleted{tout} compact $K$, \deleted{\unsure{Pour}} et toute suite $(f_i)$ dans $K^\infty_K(K)$ tendant vers une \deleted{\unsure{limite}} limite $f$ en mesure, on a $P(f)\leq\underline{\lim}\,P(f_i)$

En fait, $P_3$ est équivalente à la condition, plus forte en apparence

$P'_3)$ \inserted{Pour \deleted{tout} compact $K$,} $P$ est semi continue inférieurement sur \deleted{tout $K$} pour la topologie faible sur $L^\infty(K)$

Il faut montrer que pour \unsure{tout} réel $a$, \deleted{$P^{-1}([\unsure{a|-\infty},\unsure{-\infty|a}[)$} \inserted{l'ens $A$ des} \deleted{est \unsure{fblt} fermé} $f\in L^\infty(K)$ telles que $P(f)\leq a$ est \unsure{fblt} fermé. \deleted{D'après le th. de \unsure{Banach-Dieudonné}} \inserted{Cet ensemble est convexe en} \deleted{\unsure{on est} \illegible{2}} \deleted{il suffit de prouver que son intersection avec toute} \deleted{partie \unsure{fermée|bornée} de $L^\infty(K)$ est \unsure{fblt} \deleted{\unsure{compacte}} \unsure{fermée}.} vertu de $P_2$, \deleted{\unsure{et pour prouver}} \inserted{et en vertu} de $P_3$ il est fermé pour les suites bornées dans $L^\infty(K)$ qui convergent en mesure. Or un tel ensemble $A$ est déjà \unsure{fblt} fermé, car en vertu du th. de Banach-Dieudonné, on est ramené au cas où $A$ est borné dans $L^\infty(K)$. Il est donc \unsure{fblt} \unsure{relt} \unsure{cpt} dans $L^{\unsure{1|2}}(K)$, \deleted{\unsure{donc}} de plus il est \uline{fermé} dans $L^1(K)$ (pour la top. forte, donc pour la
% archive sheet number: 13

\scanpage{15}
topologie faible) donc \unsure{fblt} \inserted{\unsure{relt}} compact dans $L^1(K)$ ; \deleted{d'où} \inserted{i.e.} \deleted{facilement \unsure{qu'il|que il} est \unsure{fblt} compact dans $L^\infty(K)$.} pour $\sigma(\deleted{\unsure{L^1}}L^\infty(\unsure{K|k})\unsure{,}L^\infty(K))$, donc (étant borné dans $L^\infty(K)$), pour $\sigma(L^\infty(K),L^1(K))$.
% slanted
\margin{\unsure{Comme on|Comme} \unsure{peut|on|va} \unsure{supposer|suppose} $\unsure{a|\omega}$ fini}

Considérons maintenant l'ens $A_{\unsure{P|p}}(K)$ dans $\underline R\times L^\infty(K)$ formé des $(\lambda,f)$ tels que $\lambda\geq P(f)$. $P_2$ signifie que $A_P$ est convexe, $P'_3$ que $A_P$ est \unsure{fblt} fermé, donc (Hahn Banach) $A_P$ est une intersection de demi-espaces \unsure{fblt} fermés,
% slanted
\margin{fait général pour les \unsure{fonct.} convexes semi continues inférieurement} qui sont forcément de la forme $\lambda\geq L_i(f)+\deleted{\unsure{a}}a_i$ ($L_i$ \unsure{formes|forme} \unsure{fblt} \unsure{continues|continue} sur $L^\infty(K)$, i.e. $L_i\in L^1(K)$ ; \deleted{\unsure{$L_i$} \unsure{en vertu de}} $a_i\in\underline R$) d'où \deleted{$P(f)=\unsure{\mathop{\rm Sup}}$} $P(f)=\mathop{\rm Sup}\limits_i(L_i(f)+a_i)$. En vertu de $P_1$ on a \deleted{\unsure{$P(f)$}} aussi, pour \deleted{tout} \deleted{\unsure{$g$}} \inserted{$i$ et \deleted{tout} $g\in L^\infty(K)$} \unsure{spectralt|spectralement} \unsure{équivalent} à $f$, $P(f)\geq L_i(g)+a_i$ d'où \inserted{en vertu de prop. $2$,} faisant varier $g$ $P(f)\geq\int\varphi_{L_i}\varphi_f+a_i$. Posons $\varphi_{L_i}\unsure{:|=}\varphi_i$ ; on a donc
\begingroup\csname tagsleft@true\endcsname
\[
\boxed{P(f)=\mathop{\rm Sup}\limits_{i\in I}a_i+\int_0^{\unsure{\lvert K\rvert}}\varphi_i\varphi_f}
\tag*{(1)}
\]
\endgroup
% archive sheet number: 16

\scanpage{24}
\begingroup\csname tagsleft@true\endcsname
\[
\deleted{\Delta_A(t)=\Delta_{\lvert A\rvert}(t)=\Delta_{\varphi_{\lvert A\rvert}}(t)
=\exp\int_0^t\log\varphi_{\lvert A\rvert}(s)\,ds}
\tag*{\deleted{(5)}}
\]
\[
\deleted{\Delta(A)=\deleted{\unsure{\int_0^{T(1)}}}\Delta_A(\omega)
\qquad(\text{si }\{\omega=T(1)\}<+\infty)}
\tag*{\deleted{(6)}}
\]
\endgroup

\deleted{\unsure{Or|$\varphi$} Par suite, $\Delta_A(t)$ est une fonction \deleted{\unsure{finie}} continue de $t$ à valeurs $0\leq\Delta_A(t)<\infty$. $\Delta(A)$ est \deleted{\unsure{tjrs}} finie si $\unsure{\mathrm{m.}}T(1)<\infty$ (i.e. si $T$ est une trace finie), \unsure{sinon} \unsure{Si|si} la trace est \unsure{finie|bornée}, on a $1\leq\Delta(A)<+\infty$.}

% slanted
\margin{\unsure{Définition} de \deleted{\unsure{de la forme}} $\varphi_f=\unsure{1+f|1+f^+}$, $f\in\tilde b$ $(\unsure{\bar b|\tilde b}=C+\underline b)$. Définition de $\unsure{W|W^*}(A)$. les ``$A$'' \unsure{$\in$|de} l'idéal \deleted{\unsure{de finie}} $\unsure{\underline a_0|a_0}$ pour \unsure{$\equiv$}} Supposons que $\underline C$ soit une sous-algèbre commutative \unsure{auto-adjointe} maximale de $\underline b$, $\underline C=\mathcal C_0(M)$, $m$ la mesure sur $M$ définie par la \unsure{restriction de|restriction à} $T$. Si $\unsure{t|s}\in M$ est un point de $M$ de mesure non nulle, on voit \unsure{bien|par} \unsure{(en vertu de|(l'étude de} la maximalité de $\underline C$) que \inserted{$E=\unsure{U_{\varphi_{\{t\}}}|\!U\varphi_{\{t\}}}$} \deleted{\unsure{$E_\varphi$}} est un projecteur minimal de \deleted{\unsure{$\mathcal A$}} \unsure{$\underline b$} \unsure{Réciproquement} \unsure{soit|si} $E$ est un projecteur minimal dans \inserted{\unsure{$\mathcal A$}} \deleted{\unsure{$\underline E$}, $\underline C$ \inserted{$=\mathcal C_0(M)$} une algèbre \unsure{autoadjointe} maximale de $\underline b$, contenant $E$, alors \deleted{il est} \deleted{\unsure{immédiat}} \unsure{identique} à la fonction \deleted{$E$ est la fonction} caractéristique d'un \deleted{point} \deleted{ensemble à la fois} \deleted{ouvert et \unsure{fermé}} \unsure{ensemble} ouvert et fermé de $M$, \deleted{\unsure{vide} formé} \unsure{évident|réduit} à un point donc d'une \inserted{donc} mesure $>0$. \unsure{Ceci|On} dit \deleted{\unsure{qu'un} \unsure{l'opérateur}} \inserted{que la} \inserted{trace} $T$ est \deleted{\unsure{continue}} non atomique si il n'existe pas de projecteur $>0$ minimal dans $\underline{\unsure{C|b}}$, i.e. si \deleted{\unsure{on ne}} la mesure $\unsure{m|\mu}$ sur un espace $M$. i.e. de à \unsure{une} \unsure{$\underline C$} \inserted{$=\mathcal C_0(M)$} \unsure{associé est} non atomique \unsure{tout} \unsure{$C^*$} \unsure{fermé} de $\mathcal A$.} \inserted{Si $E$ \uline{est} un} projecteur $>0$ \deleted{\unsure{nul}} minimal, \deleted{\unsure{dans un}} on voit facilement qu'il existe un tel
% archive sheet number: 25

\scanpage{27}
\uline{Prop 6} \deleted{Supposons $A$ \unsure{non-atomique}} \inserted{Soit $A\in\underline b^+$, on a}
\begingroup\csname tagsleft@true\endcsname
\[
\begin{aligned}
\varphi_A^-(t)&=\mathop{\rm Inf}\limits_{\substack{E\in P(\mathcal A)\\\mathrm{Tr}\,E<t}}
\mathop{\rm Sup}\limits_{x\in(1-E).H}(Ax,x)\\
&=\mathop{\rm Inf}\limits_{\substack{E\in P(\mathcal A)\\\mathrm{Tr}\,E<t}}\lVert(1-E)A(1-E)\rVert
\end{aligned}\tag*{(5)}
\]
\endgroup

\uline{Démonstration} Soit $\underline C$ \inserted{$=\mathcal C_0(M)$} une sous-algèbre commutative autoadjointe \unsure{maximale} dans $\underline b$, contenant $A$. Si on se borne dans \unsure{le calcul des|le calcul du} $\mathop{\rm Inf}$ aux $E$ qui sont dans $\underline C$, on trouve
\[
\mathop{\rm Inf}\limits_{\substack{E\ \text{\unsure{mesurable|ensemble}}\subset M\\\mathrm{Tr}(E)<t}}
\lVert(1-E)\deleted{\unsure{E}}A\rVert_\infty\unsure{\geq|=}\varphi_A^-(t)
\]
\figure{\begin{tikzpicture}[x=0.46cm,y=0.46cm,line cap=round]
\draw (0,3.2) .. controls (.32,2.42) and (.63,1.72) .. (.9,1.64);
\draw (.67,2.15) -- (1.0,1.72);
\draw (1.1,1.63) .. controls (1.65,1.57) and (2.25,1.47) .. (2.85,1.47);
\draw (2.13,1.51) -- (2.1,.06); \fill (2.1,.06) circle (.045);
\draw (3,1.13) -- (3.56,.53);
\end{tikzpicture}}
(N\textsuperscript{o} 1, \unsure{formule} (4)) \deleted{\unsure{d'où|d'ici}} \unsure{donc} dans (5) $\varphi^-_A(t)$ est $\geq$ le deuxième mb.

% slanted
\margin{on pose $\lvert A\rvert=\sqrt{A^*A}$} \uline{Corollaire 1} $A\in\underline b$ implique
\begingroup\csname tagsleft@true\endcsname
\[
\begin{aligned}
\varphi^-_{\lvert A\rvert}(t)&=\mathop{\rm Inf}\limits_{\substack{E\in P(\mathcal A)\\\unsure{\mathrm{Tr}|\mathrm T}(E)<t}}
\mathop{\rm Sup}\limits_{x\in(1-E)H}\lVert Ax\rVert\\
&=\mathop{\rm Inf}\limits_{\substack{E\in P(\mathcal A)\\T(\unsure{E|t})<t}}\lVert A(1-E)\rVert
\end{aligned}\tag*{(6)}
\]
\endgroup

\uline{Corollaire 2} $A\in\underline b$, $B\unsure{\in|\epsilon}\mathcal A$ impliquent
\begingroup\csname tagsleft@true\endcsname
\[
\begin{aligned}
&\left\{\begin{aligned}
\varphi_{\lvert BA\rvert}&\leq\lVert B\rVert\varphi_{\lvert A\rvert}\\
\varphi_{\lvert AB\rvert}&\unsure{\leq|=}\lVert B\rVert\varphi_{\lvert A\rvert}
\end{aligned}\right.
\end{aligned}\tag*{(7)}
\]
\endgroup
% slanted
\margin{\uline{Corollaire} 5 $\varphi_{\lvert1+A\rvert}\leq\varphi_{1+\lvert A\rvert}$}

\uline{Corollaire 3} $\varphi_A$ est fonction croissante de $A\in\underline b^+$

\uline{Corollaire 4} Les fonctions $A\unsure{\text{--}|\mapsto}\varphi^-_{\lvert A\rvert}(t)$ sur $b$ forment un ensemble équicontinu de fonctions (quand $\underline b$ est \unsure{muni} de sa topologie \unsure{normée|usuelle})
% archive sheet number: 28

\scanpage{35}
\deleted{On suppose la trace continue.}

\uline{Astuce} : On plonge \unsure{un|une} \unsure{anneau|algèbre} avec une trace non continue, dans un \unsure{anneau} avec une trace continue.

Si \deleted{\unsure{$E_{\varphi}$} \unsure{$\int_{T(A)}$} \unsure{$(t)=$}}

\uline{Prop} 1
\[
\varphi_A(t+0)=\mathop{\rm Inf}\limits_{\substack{E\in P\\\mathrm{Tr}\,E<t}}
\mathop{\rm Sup}\limits_{x\in(1\unsure{+|-}E)H}(Ax,x\unsure{]|)}
\qquad\text{(trace continue)}
\]

% margin note written sideways
\margin{\unsure{:} \unsure{preuve|prouver} \deleted{$\unsure{\Delta_t}(1+A)\unsure{=|\leq}\Delta_t(1+\lvert A\rvert)$} \inserted{$\varphi_{\lvert1+A\rvert}\leq \varphi_{1+\lvert A\rvert}$}} On a en effet, comme on voit aussitôt, (décomposition spectrale) $\varphi_A(t)\geq$ second mb \deleted{\unsure{en}}, Il suffit donc de prouver $\mathop{\rm Sup}\limits_{x\in(1-E)H}(Ax,x)\geq\varphi^-_A(t)$ \inserted{si $\mathrm{Tr}\,E<t$} On est ramené au cas où $A$ est égal à $\lambda F$, $F$ projecteur de trace $\unsure{\geq|>}t$, donc \unsure{à|si} $t$ \unsure{près|fini} \unsure{$=$}. \deleted{\unsure{car ceci}} \inserted{(car $A\geq\varphi_A^{\unsure{-}}(t)F$ avec $\mathrm{Tr}\,F\geq t$)} Il suffit de prouver : Si \uline{$\mathrm{Tr}\,E<\mathrm{Tr}\,F$, alors $(1-E)H\cap FH\neq(0)$}.

En effet, on a \deleted{\unsure{$E$}}, pour
\[
\begin{aligned}
\mathcal X&=E\mathcal H,\ \mathcal Y=F\mathcal H,\\
\mathcal X'&=\mathcal H\ominus\mathcal X,\ \mathcal Y'=\unsure{\mathcal H|\mathcal X}\ominus\mathcal Y :
\end{aligned}
\]
$\mathcal X\ominus(\mathcal X\cap\mathcal Y')\sim\mathcal Y\ominus(\mathcal Y\cap\mathcal X')$ donc $\mathcal Y\cap\mathcal X'=0$ impliquerait \deleted{$\mathcal X\unsure{\sim|=}\unsure{0|\ominus}$} $\mathcal Y\unsure{\prec|\lesssim}\mathcal X$ d'où $\mathrm{Tr}\,F\leq\mathrm{Tr}\,E$, absurde.

\uline{Corollaire} 1 $A$ quelconque implique
\[
\varphi_{\lvert A\rvert}(t+0)=\mathop{\rm Inf}\limits_{\substack{E\in P\\\mathrm{Tr}\,E<t}}
\mathop{\rm Sup}\limits_{x\in(1-E)H}\lVert Ax\rVert\qquad\text{(trace continue)}
\]

\uline{Corollaire 2}
\[
\begin{aligned}
\varphi_{\lvert PA\rvert}\deleted{\unsure{(t)}}&\leq\lVert P\rVert\,\varphi_{\lvert A\rvert}\deleted{\unsure{(t)}}\qquad\text{(trace quelconque),}\\
\varphi_{\lvert AQ\rvert}&\leq\lVert Q\rVert\,\varphi_{\lvert A\rvert}\qquad\deleted{\unsure{\varphi_{\lvert1+A\rvert}\leq\varphi_{1+\lvert A\rvert}}}
\end{aligned}
\]

\uline{Corollaire 3} $\varphi_A$ croît avec \unsure{l'hermitien|l'élément} positif \deleted{\unsure{$A$ positif}} \unsure{$A$}

% margin note written sideways
\margin{\uline{Corollaire 4} $\varphi_{\lvert1+A\rvert}\leq\varphi_{1+\lvert A\rvert}=1+\varphi_{\lvert A\rvert}$}
% slanted
\margin{$\varphi_{\alpha(A)}=\alpha\varphi_A$} \uline{Corollaire 5} Équicontinuité \inserted{des fonctions $A\to\varphi_{\lvert A\rvert}(t)$} \deleted{\unsure{des} $A$ $\varphi_{\lvert A\rvert}(t\unsure{+0})$ en fonction de $A$.}

\uline{Prop} 2 $\varphi_{\lvert A\rvert}=\varphi_{\lvert A^*\rvert}$ Il suffit de le prouver \deleted{que} $\varphi_{AA^*}=\varphi_{A^*A}$ et de le prouver pour $A\in\unsure{a|\mathfrak a}$. \inserted{Posons} \deleted{$A\in\underline a$, \unsure{et} \illegible{1} le prouver} $\unsure{K|X}_{\lvert\unsure{K|A}\rvert A} = K_{\lvert A\rvert}\unsure{\theta\theta|\theta^*}$ ou \unsure{$\int$} \deleted{\unsure{$x\,d.K_{\lvert A\rvert}(xp)|x\,d.K_{\lvert A\rvert}(x,p)$}}
% archive sheet number: 36

\scanpage{39}
\deleted{Rappelons la \unsure{caractérisation|construction}}

Notations $(\lambda_i(u))$ si $u\in L_0(E)$ : c'est la suite des valeurs \deleted{\unsure{caractéristiques}} propres \inserted{non nulles} de $u$, rangées par ordre de modules \unsure{non} \unsure{décroissants|croissants}, complétée par des $0$ si \unsure{cette} suite est finie.

Si la dim de $E$ est finie, \inserted{soit $n$} il y a lieu \inserted{d'attacher} \margin{?} son intérêt \deleted{\unsure{seulement}} à $\lambda_1\unsure{(u)},\ \unsure{\ldots}\ \lambda_n(u)$ \inserted{Si le rang est fini, \unsure{on s'arrête|on garde} même \deleted{\unsure{au terme|au rang}} à partir du rang $\unsure{n|r}+1$} On pose de plus \deleted{$\unsure{\rho|f}_n(u)\unsure{=|:}$} , pour $u\in L_0(E,F)$
\[
\boxed{\boxed{\rho_n(u)=\lambda_n(\lvert u\rvert)}}
\]
donc $\rho_n(u)$ est une \unsure{suite} \unsure{décroiss\textsuperscript{te}|décroissante|décroist} de nbs positifs.

\[
\deleted{\unsure{\rho_1(u)\leq\lVert u\rVert}}\qquad\deleted{\unsure{\lambda_1(u)}}
\]
\[
\boxed{\lvert\lambda_1(u)\rvert\leq\rho_1(u)=\lVert u\rVert}
\qquad\boxed{\rho_n(u)=\rho_n(u^*)}
\]

Si $h$ est \deleted{\unsure{hermitien}} positif, on sait que
\[
\boxed{\rho_n(h)=\mathop{\rm Inf}\limits_{E_{n-1}}
\deleted{\mathop{\rm Sup}\limits_{\unsure{x\in E}}\lVert\unsure{hx}\rVert}
\unsure{(}\mathop{\rm Sup}\limits_{\substack{x\in E_{n-1}\\\lVert x\rVert\leq1}}(hx,x)\bigr)\deleted{{}^{\unsure{1/2|2}}}}
\]
% slanted
\margin{$\unsure{\mathscr H_0|\mathscr H_\infty}$} où \deleted{$h_F$ est la restriction à $F\times F$ de la forme} \inserted{$E_{n-1}$ parcourt les} sous-espaces vectoriels fermés de codimension $\unsure{n-1|n+1|n}$. \deleted{sesquilinéaire $(hx,y)$}

On en conclut $\rho_n(h^{\unsure{4|2|a}})^2=\rho_n(h^{\unsure{2|7}})=\mathop{\rm Inf}\limits_{F_n}\mathop{\rm Sup}\limits_{\substack{x\in F_n\\\lVert x\rVert\leq1}}\lVert hx\rVert^{\unsure{2|4}}$, \deleted{\unsure{d'où}} d'où
\[
\begin{aligned}
&\deleted{\rho_n(h)=\mathop{\rm Inf}\mathop{\rm Sup}\limits_{x\in F_n}hx}\\
&\deleted{\text{d'où }\rho_n(u)^2=\rho_n(u^*u)=\mathop{\rm Inf}\limits_{\unsure{E_n}}\mathop{\rm Sup}\limits_{x\in E_n}\lVert u^*ux\rVert}\\
&\rho_n(h)\unsure{'}=\mathop{\rm Inf}\limits_{E_n}\mathop{\rm Sup}\limits_{\substack{x\in E_{n-1}\\\lVert x\rVert\leq1}}\lVert hx\rVert\deleted{{}^{\unsure{2}}}\qquad \deleted{\unsure{(h\geq0)}}\ (h\geq0)
\end{aligned}
\]
D'où, en posant $h=\lvert u\rvert$
\[
\boxed{\rho_n(u)=\mathop{\rm Inf}\limits_{F_n}\mathop{\rm Sup}\limits_{\substack{x\in E_{n-1}\\\lVert x\rVert\leq1}}\lVert ux\rVert}
\]

On \deleted{\unsure{en conclut}} \inserted{en \unsure{conclut}} aussitôt :
\[
\boxed{\begin{aligned}
\rho_n(uv)&\leq\lVert u\rVert\,\rho_n(v)\\
\rho_n(uv)&\leq\lVert v\rVert\,\rho_n(u)
\end{aligned}}
\]
Et de \unsure{m\textsuperscript{me}|même} : \uline{$\rho_n(h)$ est fonction croissante de l'hermitien positif $h$.} \deleted{\uline{Corollaire} Soit $u\in L_0(E,F)$} \unsure{Enfin|Écrivons} \deleted{\unsure{on}}
\[
\deleted{\boxed{\unsure{\sum_i}\rho_i(u)=\mathop{\rm Sup}\limits_{\unsure{p_n}}\unsure{\mathrm{Tr}}\,\lvert\unsure{p_nu}\rvert}}
\]

\margin{?} \deleted{On \unsure{en} \unsure{démontre} \unsure{Tr} \illegible{1}} les inégalités de Ky-Fan
\[
\boxed{\begin{aligned}
\rho_{m\unsure{+|-}n\unsure{-|+}1}(A\unsure{B|\Pi})&\unsure{\leq|=}\rho_{m+1}(A)\rho_{n+1}(B)\\
\rho_{m\unsure{+|-}n\unsure{+|-}1}(A+\unsure{B|\Pi})&\unsure{<|\leq}\rho_{m+1}(A)+\rho_{n+1}(B)
\end{aligned}}
\]
% archive sheet number: 40

\scanpage{49}

séparé $E\otimes F$, \deleted{\unsure{dont|avec|dans}}\inserted{où} le produit scalaire est défini par (3). La norme d'un élément \deleted{\unsure{est notée}} $u\in E\overset{(2)}{\otimes}F$ est notée $\lVert u\rVert_2$\unsure{.} $E\overset{(2)}{\otimes}F$ donc un espace de Hilbert.

\uline{Prop} 1 \unsure{Soient} $E$ et $F$ deux espaces de Hilbert. \deleted{\unsure{1°} \unsure{Si} $x\in E$, $y\in F$, on a $\lVert x\otimes y\rVert_2=\lVert x\rVert\lVert y\rVert$. \unsure{2°}} Soit $(e_i)_{i\in I}$ une base orthonormale dans $E$, $(\unsure{b_j|f_j|b})_{j\in J}$ une base orthonormale dans $F$, alors $(e_i\otimes b_j)_{i,j\in I\times J}$ est une base orthonormale dans $E\overset{(2)}{\otimes}F$.

En effet, d'après (4) le système $(e_i\otimes \unsure{f_j|b_j})$ est orthonormal; \deleted{\unsure{donc pour la totalité} \illegible{2}} \deleted{\unsure{montrons} qu'il est \unsure{dense} dans $E\overset{(2)}{\otimes}F$.} \deleted{\unsure{désignons}} \deleted{En effet, \illegible{2}} \inserted{\unsure{montrons}} \deleted{\unsure{il nous reste à}} que l'espace vectoriel qu'il engendre est dense dans $E\overset{(2)}{\otimes}F$. \deleted{Pour ceci,} Si $E_0$ (resp $F_0$) \deleted{est}\inserted{désigne} l'espace engendré par les $e_i$ (resp. les $f_j$) \deleted{il suffit de} \deleted{\unsure{est évident en effet,}} \deleted{\unsure{fait que|faut que}} $E_0$ est dense dans $E$, $F_0$ dense dans $F$, donc $E_0\otimes F_0$ dense dans $E\otimes F$, donc aussi dans $E\overset{(2)}{\otimes}F$.

Notons \uline{aussi la formule}
\begingroup
\csname tagsleft@true\endcsname
\[
\boxed{\lVert x\otimes y\rVert_2=\lVert x\rVert\lVert y\rVert}
\tag*{(\unsure{7})}
\]
\endgroup
équivalente en effet à \deleted{\unsure{la} \illegible{1}} \deleted{qui résulte \unsure{immédiatement de}} $(x\otimes y,x\otimes y)=(x,x)(y,y)$ \deleted{\unsure{l'égalité qui précède} \illegible{2}}

Notons aussi que l'isomorphisme \inserted{canonique} $\overline{E\otimes F}=\overline{E}\otimes\overline{F}$ (formule (1)) conserve \unsure{manifestement|manifest\textsuperscript{t}} les structures préhilbertiennes, d'où un isomorphisme canonique
\begingroup
\csname tagsleft@true\endcsname
\[
\overline{E\overset{(2)}{\otimes}F}=\overline E\overset{\unsure{(1)|(2)}}{\otimes}\overline F
\tag*{\unsure{(8)|8)}}
\]
\endgroup
% archive sheet number: 50

